\exercisesection{\S~1}

\begin{enumerate}
\def\labelenumi{\arabic{enumi})}
\tightlist
\item
  设 \(\mathrm{G}\) 是一个有限维、连通、交换的复李群，且 \(\mathrm{V}\) 是其李代数。
\end{enumerate}

\begin{enumerate}
\def\labelenumi{\alph{enumi})}
\item
  映射 \(\exp_{\mathrm{G}}: \mathrm{V} \to \mathrm{G}\) 是满同态；其核 \(\Gamma\) 是 \(\mathrm{V}\) 的离散子群。
\item
  G 紧当且仅当 \(\Gamma\) 是 V 中的一个格；此时称 G 为复环面。
\item
  设 \(\Gamma\) 是 \(\mathbf{C}^2\) 的离散子群，由元素 \(e_1 = (1,0), e_2 = (0,1), e_3 = (\sqrt{2},i)\) 生成；令 \(\mathrm{G} = \mathrm{C}^2 / \Gamma\)，\(\mathrm{H} = (\Gamma + \mathrm{C}e_1) / \Gamma\)。证明 H 同构于 \(\mathbf{C}^*\)，G/H 是维数为 1 的复环面，但 G 不包含非零复环面。
\item
  每个有限维连通紧复李群都是复环面（参见~第 III 章，§6，节 3，命题 6）。
\end{enumerate}

\begin{enumerate}
\def\labelenumi{\arabic{enumi})}
\setcounter{enumi}{1}
\tightlist
\item
  设 H 是由复数 \(\tau\) 中满足 \(\mathcal{I}(\tau) > 0\) 者组成的集合。
\end{enumerate}

\begin{enumerate}
\def\labelenumi{\alph{enumi})}
\item
  证明可以在 H 上定义离散群 \(\mathbf{SL}(2,\mathbf{Z})\) 的解析左作用律，规定 \(\gamma\tau=\frac{a\tau+b}{c\tau+d}\)；对所有 \(\gamma=\begin{pmatrix}a&b\\ c&d\end{pmatrix}\in\mathbf{SL}(2,\mathbf{Z})\) 和 \(\tau\in H\)，该规定成立。
\item
  对于 \(\tau\in H\)，记 \(T_{\tau}\) 为复环面 \(\mathbf{C}/(\mathbf{Z}+\mathbf{Z}\tau)\)。证明映射 \(\tau\mapsto T_{\tau}\) 通过取商诱导出从集合 \(\mathrm{H}/\mathrm{SL}(2,\mathbf{Z})\) 到维数为 1 的复环面的同构类集合的一个双射。
\end{enumerate}

\begin{enumerate}
\def\labelenumi{\arabic{enumi})}
\setcounter{enumi}{2}
\item
  设 G 是 \(\mathbf{O}(n)\) 的一个整子群，使得 G 在 \(\mathbf{R}^{n}\) 上的恒等表示是不可约的。证明 G 是闭的（将 G 写成 \(K \times N\) 的形式，其中 K 是紧的、N 是交换的，并证明 \(\dim \overline{N} \leq 1\)）。
\item
  证明命题 3（节 3）中的条件分别等价于下列每个条件：
\end{enumerate}

(ii') 群 \(\operatorname{Ad}(\mathbf{G})\) 在 \(\operatorname{Aut}(\mathbf{L}(\mathbf{G}))\) 中相对紧。

(ii'\,') 群 \(\operatorname{Ad}(\mathbf{G})\) 在 \(\operatorname{End}(\mathbf{L}(\mathbf{G}))\) 中相对紧。

\begin{enumerate}
\def\labelenumi{(\alph{enumi})}
\setcounter{enumi}{21}
\tightlist
\item
  G 的单位元 e 的每个邻域都包含一个在内自同构下稳定的 e 的邻域（参见~《积分》，§3，节 1，命题 1）。
\end{enumerate}

\begin{enumerate}
\def\labelenumi{\arabic{enumi})}
\setcounter{enumi}{4}
\tightlist
\item
  设 A 是 \(\mathbf{GL}(3,\mathbf{R})\) 的闭子群，由矩阵 \((a_{ij})\) 组成，这些矩阵满足 \(a_{ij}=0\)（当 i\textgreater j 时）、\(a_{ii}=1\)（\(1\leq i\leq3\)）、\(a_{12}\in \mathbf{Z}\)、\(a_{23}\in \mathbf{Z}\)；令 \(\theta\in \mathbf{R}-\mathbf{Q}\)，并令 B 是 A 的子群，由矩阵 \((a_{ij})\) 组成，这些矩阵满足 \(a_{12}=a_{23}=0\) 且 \(a_{13}\in\theta\mathbf{Z}\)，并令 G=A/B。
\end{enumerate}

\begin{enumerate}
\def\labelenumi{\alph{enumi})}
\item
  证明 G 是一个维数为 1 的李群，且其李代数是紧的。
\item
  证明 \(\mathrm{C}(\mathrm{G}) = \overline{\mathrm{D}(\mathrm{G})} = \mathrm{G}_{0}\)，并且 \(G_{0}\) 是 G 的最大紧子群。
\item
  证明 G 不是 \(G_{0}\) 与任何子群的半直积。
\end{enumerate}

\begin{enumerate}
\def\labelenumi{\arabic{enumi})}
\setcounter{enumi}{5}
\item
  证明一个（实）李代数是紧的，当且仅当它存在一组基 \((e_{\lambda})_{\lambda\in\mathrm{L}}\)，使得结构常数 \(\gamma_{\lambda\mu\nu}\) 关于 \(\lambda,\mu,\nu\) 反对称（即它们满足 \(\gamma_{\lambda\mu\nu}=-\gamma_{\mu\lambda\nu}=-\gamma_{\lambda\nu\mu}\)）。
\item
  对合李代数是指一个配备自同构 s 的（实）李代数 \(\mathfrak{g}\)，满足 \(s \circ s = 1_{\mathfrak{g}}\)。记 \(\mathfrak{g}^{+}\)（分别为 \(\mathfrak{g}^{-}\)）为 \(\mathfrak{g}\) 中 s 对应于特征值 +1（分别为 -1）的特征子空间。
\end{enumerate}

\begin{enumerate}
\def\labelenumi{\alph{enumi})}
\item
  证明 \(\mathfrak{g}^{+}\) 是 \(\mathfrak{g}\) 的子代数，且 \(\mathfrak{g}^{-}\) 是一个 \(\mathfrak{g}^{+}\)-模；有 \([\mathfrak{g}^{-},\mathfrak{g}^{-}] \subset \mathfrak{g}^{+}\)，并且 \(\mathfrak{g}^{+}\) 与 \(\mathfrak{g}^{-}\) 关于 Killing 型正交。
\item
  证明下列条件等价：
\end{enumerate}

\begin{enumerate}
\def\labelenumi{(\roman{enumi})}
\item
  \(\mathfrak{g}^{+}\)-模 \(\mathfrak{g}^{-}\) 是单模。
\item
  代数 \(\mathfrak{g}^{+}\) 在所有不同于 \(\mathfrak{g}\) 的 \(\mathfrak{g}\) 子代数中是极大的。若这些条件成立，且 \(\mathfrak{g}^{+}\) 不包含 \(\mathfrak{g}\) 的非零理想，则称对合李代数 \((\mathfrak{g}, s)\) 为不可约的。
\end{enumerate}

\begin{enumerate}
\def\labelenumi{\alph{enumi})}
\setcounter{enumi}{2}
\tightlist
\item
  假设 \(\mathfrak{g}\) 是半单的。证明下列条件等价：
\end{enumerate}

\begin{enumerate}
\def\labelenumi{(\roman{enumi})}
\item
  \(\mathfrak{g}\) 的在 \(s\) 下稳定的理想只有 \(\{0\}\) 和 \(\mathfrak{g}\)；
\item
  \(\mathfrak{g}\) 或者是单的，或者是两个在 \(s\) 下互换的单理想之和。
\end{enumerate}

证明当 \((\mathfrak{g}, s)\) 不可约时这些条件成立（注意，\(\mathfrak{g}\) 是满足 (ii) 且在 s 下稳定的理想的直和）。

\begin{enumerate}
\def\labelenumi{\alph{enumi})}
\setcounter{enumi}{3}
\item
  从现在起假设 \(\mathfrak{g}\) 是紧半单的。证明对合李代数 \((\mathfrak{g}, s)\) 不可约，当且仅当 s 不等于恒等映射，且 \((\mathfrak{g}, s)\) 满足 c) 中的等价条件（令 p 为 \(\mathfrak{g}^{+}\)-模 \(\mathfrak{g}^{-}\) 的一个子模，q 为其关于 Killing 型的正交补；注意到 \([p, q] = 0\)，并由此推出 \(p + [p, p]\) 是 \(\mathfrak{g}\) 的理想）。
\item
  证明 \(\mathfrak{g}\) 是一族在 s 下稳定的理想 \((\mathfrak{g}_{i})_{0\leq i\leq n}\) 的直和，其中 \(\mathfrak{g}_{0}\) 由 s 固定，且对合代数 \((\mathfrak{g}_{i},s|\mathfrak{g}_{i})\) 在 \(1\leq i\leq n\) 时不可约。
\end{enumerate}

\begin{enumerate}
\def\labelenumi{\arabic{enumi})}
\setcounter{enumi}{7}
\tightlist
\item
  设 G 是一个紧半单李群，u 是 G 的一个阶为 2 的自同构，K 是 u 的不动点集的单位元连通分支，X 是流形 G/K（齐性空间 X 因而称为对称空间）。
\end{enumerate}

\begin{enumerate}
\def\labelenumi{\alph{enumi})}
\item
  证明，如果 G 是概单的，则 K 在所有不同于 G 的连通闭子群中是极大的；换言之（第 III 章，§3，习题 8），G 在 X 上的作用是本原的。对称空间 X 因而称为不可约的。
\item
  假设流形 X 是单连通的；证明它同构于若干流形的乘积，其中每个流形都同构于一个李群或一个不可约对称空间。
\end{enumerate}

\begin{enumerate}
\def\labelenumi{\arabic{enumi})}
\setcounter{enumi}{8}
\tightlist
\item
  设 \(\mathfrak{a}\) 是一个实或复李代数，且 \(G\) 是 \(\operatorname{Aut}(\mathfrak{a})\) 的一个紧子群。证明 \(\mathfrak{a}\) 有一个在 \(G\) 下稳定的 Levi 子代数（第 I 章，§6，节 8，定义 7）（归结到 \(\mathfrak{a}\) 的根基为交换的情形，并使用《积分》，第 VII 章，§3，节 2，引理 2）。
\end{enumerate}

